\documentclass[11pt]{article}

\usepackage[margin=1in]{geometry}
\usepackage{amsmath, amssymb}
\usepackage{graphicx}
\usepackage{booktabs}
\usepackage{xcolor}
\usepackage{listings}
\usepackage{hyperref}

\lstset{
  language=Python,
  basicstyle=\ttfamily\small,
  keywordstyle=\color{blue!70!black},
  commentstyle=\color{green!40!black},
  stringstyle=\color{red!60!black},
  frame=single,
  breaklines=true,
  showstringspaces=false,
  columns=fullflexible,
}

\title{MATH 7243 --- Lab 2: Methods in Linear Regression}
\author{Wilson Glass}
\date{October 3, 2026}

\begin{document}
\maketitle

%=============================================================================
\section*{Problem 1: Bootstrapping a Confidence Interval}

\subsection*{Setup}

We use the Ames housing data, cleaned as in the Master notebook: remove the outliers
with \texttt{GrLivArea} + \texttt{BsmtUnfSF} $\ge 4000$, keep only the numerical
columns, and drop the categorical \texttt{MSSubClass}. This leaves $n = 1436$ houses.
We fit the one-variable model
\[
  \texttt{SalePrice} = \beta_0 + \beta_1 \cdot \texttt{1stFlrSF}.
\]

\begin{lstlisting}
import io
import requests
import numpy as np
import pandas as pd
from matplotlib import pyplot as plt

url = "https://raw.githubusercontent.com/tipthederiver/Math-7243-2020/master/Datasets/Ames/train.csv"
ames = pd.read_csv(io.StringIO(requests.get(url).content.decode('utf-8')))

z = ames['GrLivArea'] + ames['BsmtUnfSF'] < 4000
data = ames[z]
data = data.select_dtypes(include=['int64', 'float64'])
data = data.drop(columns='MSSubClass')     # 1436 rows
\end{lstlisting}

To bootstrap, we draw a sample of $N = 1436$ rows from the data \emph{with
replacement}, fit $\beta = (X^TX)^{-1}X^Ty$ on that sample, and store $\beta_0$ and
$\beta_1$. Repeating this 1000 times gives 1000 estimates of each parameter.

\begin{lstlisting}
N = 1000
beta0 = np.zeros(N)
beta1 = np.zeros(N)

for i in range(N):
    # 1. resample rows of the data, with replacement
    sample = data.sample(n=len(data), replace=True)

    # 2. build X (with a column of 1's) and Y
    X = np.matrix(sample["1stFlrSF"]).reshape(-1, 1)
    Y = np.matrix(sample["SalePrice"]).reshape(-1, 1)
    Xa = np.append(np.ones(X.shape), X, 1)

    # 3. fit with the normal equation
    betas = (Xa.T*Xa).I*Xa.T*Y

    # 4. store
    beta0[i] = betas[0, 0]
    beta1[i] = betas[1, 0]
\end{lstlisting}

%-----------------------------------------------------------------------------
\subsection*{Part 1: Histograms of $\beta_0$ and $\beta_1$}

\begin{lstlisting}
fig, axes = plt.subplots(1, 2)
fig.set_size_inches(12, 4)

axes[0].hist(beta0, bins=30)
axes[0].set_title("Bootstrap distribution of $\\beta_0$ (intercept)")
axes[0].set_xlabel("$\\beta_0$"); axes[0].set_ylabel("Count")

axes[1].hist(beta1, bins=30)
axes[1].set_title("Bootstrap distribution of $\\beta_1$ (slope)")
axes[1].set_xlabel("$\\beta_1$"); axes[1].set_ylabel("Count")

plt.show()
\end{lstlisting}

\begin{figure}[h]
  \centering
  \includegraphics[width=\textwidth]{figures/bootstrap_hist.png}
  \caption{Bootstrap distributions of $\beta_0$ and $\beta_1$ over 1000 resamples.
  The red dashed lines mark the bootstrap 95\% confidence interval (Part 2).}
\end{figure}

Both distributions are roughly bell-shaped, centered near the full-data estimates
$\hat\beta_0 \approx 37{,}566$ and $\hat\beta_1 \approx 121.8$. Their spread shows how much
the fitted line would change if we had drawn a different sample of houses.

%-----------------------------------------------------------------------------
\subsection*{Part 2: Bootstrap 95\% Confidence Interval}

After sorting, the middle 950 of the 1000 values run from index 25 to index 974.

\begin{lstlisting}
beta0.sort()
beta1.sort()

print("95% bootstrap CI for beta0:", beta0[25], "to", beta0[974])
print("95% bootstrap CI for beta1:", beta1[25], "to", beta1[974])
\end{lstlisting}

\textbf{Answer:}
\begin{align*}
  \beta_0 &\in [\,26{,}000.33,\; 49{,}953.89\,] \\
  \beta_1 &\in [\,110.09,\; 132.32\,]
\end{align*}

%-----------------------------------------------------------------------------
\subsection*{Part 3: Confidence Interval from the Formula}

From Section 4 of the notes, the covariance matrix of the least-squares estimate
$\hat\theta = (\hat\beta_0, \hat\beta_1)$ is
\[
  \operatorname{Cov}(\hat\theta) = \hat\sigma^2 (X^TX)^{-1},
  \qquad
  \hat\sigma^2 = \frac{\text{RSS}}{n-2},
\]
where $\hat\sigma^2$ estimates the variance of the noise around the regression line, and
we divide by $n-2$ because two parameters were fit. The diagonal entries of
$\operatorname{Cov}(\hat\theta)$ are $\operatorname{Var}(\hat\beta_0)$ and
$\operatorname{Var}(\hat\beta_1)$. Their square roots are the standard errors, and the 95\%
confidence interval is $\hat\beta_j \pm 1.96\,\text{SE}(\hat\beta_j)$. Here we use all
1436 rows.

\begin{lstlisting}
X = np.matrix(data["1stFlrSF"]).reshape(-1, 1)
Y = np.matrix(data["SalePrice"]).reshape(-1, 1)
Xa = np.append(np.ones(X.shape), X, 1)
n = len(Y)

betas = (Xa.T*Xa).I*Xa.T*Y
RSS = ((Y - Xa*betas).T*(Y - Xa*betas))[0, 0]
sigma2 = RSS/(n - 2)                    # variance of the noise

var_betas = sigma2*(Xa.T*Xa).I          # Cov(theta) = variance * (X^T X)^-1
se = np.sqrt(np.diag(var_betas))        # standard errors

print("Formula 95% CI for beta0:", betas[0, 0] - 1.96*se[0], "to", betas[0, 0] + 1.96*se[0])
print("Formula 95% CI for beta1:", betas[1, 0] - 1.96*se[1], "to", betas[1, 0] + 1.96*se[1])
print("Bootstrap std of beta0:", np.std(beta0), " formula SE:", se[0])
print("Bootstrap std of beta1:", np.std(beta1), " formula SE:", se[1])
\end{lstlisting}

\textbf{Answer:}
\begin{align*}
  \beta_0 &\in [\,27{,}673.93,\; 47{,}457.58\,] \\
  \beta_1 &\in [\,113.52,\; 129.99\,]
\end{align*}

\subsubsection*{Comparison}

\begin{center}
\begin{tabular}{lcccc}
\toprule
 & \multicolumn{2}{c}{95\% CI} & \multicolumn{2}{c}{Standard error} \\
\cmidrule(lr){2-3} \cmidrule(lr){4-5}
 & Formula & Bootstrap & Formula & Bootstrap \\
\midrule
$\beta_0$ & $[27{,}674,\ 47{,}458]$ & $[26{,}000,\ 49{,}954]$ & $5{,}047$ & $6{,}088$ \\
$\beta_1$ & $[113.52,\ 129.99]$     & $[110.09,\ 132.32]$     & $4.20$    & $5.74$    \\
\bottomrule
\end{tabular}
\end{center}

The two methods agree on the center of each interval, but the bootstrap intervals
are wider: about 21\% wider for $\beta_0$ and 35\% wider for $\beta_1$. The formula
assumes the noise has the same variance $\sigma^2$ for every house. That assumption
does not hold here: prices of large houses vary much more than prices of small ones,
so the scatter of \texttt{SalePrice} against \texttt{1stFlrSF} fans out as
\texttt{1stFlrSF} increases. The formula therefore underestimates the uncertainty in the
fit. The bootstrap makes no assumption about the noise. It resamples the actual data,
so it captures this extra variability, and its intervals are the more trustworthy of
the two for this data.

%=============================================================================
\clearpage
\section*{Problem 2: Linear Methods on High-Dimensional Data}

\subsection*{Setup}

The MRI dataset contains 609 brain slice images, each $176 \times 176 = 30{,}976$
pixels. Each image is converted to grayscale (averaging the color channels), flattened
into a single row, and downsampled by keeping every 8th pixel. This gives a data matrix of
shape $609 \times 3872$. The \textbf{features} are the 3872 pixel intensities and the
\textbf{target} is the normalized whole-brain volume \texttt{nWBV} from
\texttt{labels.csv}. The other columns of \texttt{labels.csv} (age, eTIV, ASF, \dots) are
not used, since the goal is to predict \texttt{nWBV} from the images alone.

\begin{lstlisting}
import matplotlib
from sklearn.model_selection import train_test_split
from sklearn.metrics import r2_score

file_dir = 'data/'
labels = pd.read_csv(file_dir + 'labels.csv')

DS = 8                          # downsample rate
im_size = int(30976/DS)         # 3872 features

data = np.zeros([609, im_size])
for i, file_name in enumerate(labels.Filename):
    img = np.mean(matplotlib.image.imread(file_dir + file_name), axis=2).reshape(-1)
    data[i, :] = img[::DS]

X = data
y = labels["nWBV"]

x_train, x_test, y_train, y_test = train_test_split(
    X, y, test_size=.2, random_state=255
)
\end{lstlisting}

The 80/20 split gives 487 training images and 122 test images. With 3872 features
and only 487 training samples, there are more unknowns than equations, so ordinary least
squares can fit the training data perfectly and generalizes poorly. This is why
regularization is needed. Every model below is scored by its $r^2$ on the
\emph{test} set. Training $r^2$ always improves as the penalty goes to zero, so it cannot
detect overfitting.

%-----------------------------------------------------------------------------
\subsection*{Part 1: Best $\alpha$ for Ridge Regression}

We use statsmodels' \texttt{fit\_regularized} with \texttt{L1\_wt=0} (pure ridge),
which minimizes
\[
  \frac{1}{2n}\,\text{RSS} + \frac{\alpha}{2}\sum_i \beta_i^2 .
\]
Because many background pixels are constant, \texttt{add\_constant} needs
\texttt{has\_constant="add"}. Otherwise it skips adding the intercept column on the
test set.

The search was done in two stages. A coarse sweep over
$\alpha = 10^{-6}, 10^{-5}, \dots, 10^{3}$ (\texttt{np.logspace(-6, 3, 10)}) showed the
test $r^2$ peaking near $\alpha = 10^{-2}$:

\begin{center}
\begin{tabular}{lccccc}
\toprule
$\alpha$ & $10^{-3}$ & $10^{-2}$ & $10^{-1}$ & $10^{0}$ & $10^{1}$ \\
\midrule
test $r^2$ & $0.284$ & $0.398$ & $0.285$ & $0.089$ & $-0.247$ \\
\bottomrule
\end{tabular}
\end{center}

A fine sweep of 50 values between $0.001$ and $0.05$ then located the maximum:

\begin{lstlisting}
import statsmodels.api as sm

alphas = np.linspace(0.001, 0.05, 50)
r_sqds = []

for alp in alphas:
    ols = sm.OLS(y_train, sm.add_constant(x_train, has_constant="add"))
    ols_result = ols.fit_regularized(L1_wt=0, alpha=alp)
    pred = ols_result.predict(sm.add_constant(x_test, has_constant="add"))
    r_sqds.append(r2_score(y_test, pred))

plt.plot(alphas, r_sqds, "o-")
plt.xlabel("alpha"); plt.ylabel("test $r^2$")
plt.show()

best = np.argmax(r_sqds)
print("Best alpha:", alphas[best], "with test r^2 =", r_sqds[best])
\end{lstlisting}

\begin{figure}[h]
  \centering
  \includegraphics[width=0.6\textwidth]{figures/ridge_zoom.png}
  \caption{Ridge regression: test $r^2$ against $\alpha$ (fine sweep).}
\end{figure}

\textbf{Answer:} The best ridge penalty is $\boldsymbol{\alpha \approx 0.009}$, with
test $r^2 \approx 0.398$.

\textbf{Justification.} The test $r^2$ curve rises to a single peak and then falls,
which is the bias--variance tradeoff:
\begin{itemize}
  \item \textbf{$\alpha$ too small:} the penalty is too weak to stop the model from
  fitting noise among 3872 pixel coefficients with only 487 training images. Test $r^2$
  is $0.284$ at $\alpha = 0.001$ and only $0.063$ at $\alpha = 0$.
  \item \textbf{$\alpha$ too large:} the penalty shrinks all coefficients so far toward
  zero that the model can no longer capture the real relationship. Test $r^2$ falls
  steadily past the peak, to $0.336$ at $\alpha = 0.05$, and becomes negative by
  $\alpha = 10$. A negative $r^2$ means the model does worse than predicting the mean
  \texttt{nWBV} for every image.
  \item \textbf{$\alpha \approx 0.009$} balances the two and gives the highest test
  $r^2$.
\end{itemize}
The peak is fairly flat. Every $\alpha$ from about $0.006$ to $0.016$ gives test
$r^2 \ge 0.39$, so the best value is better described as ``about $0.01$'' than as exactly
$0.009$. This $\alpha$ is on statsmodels' scale. sklearn's \texttt{Ridge} drops the
$\frac{1}{2n}$ factor on the RSS, so its optimal $\alpha$ would be numerically different.

%-----------------------------------------------------------------------------
\clearpage
\subsection*{Part 2: Best $\lambda$ for Lasso Regression}

\textbf{Notation.} Here $\lambda$ is the lasso penalty strength, as in the lasso loss
\[
  \frac{1}{2n}\,\text{RSS} + \lambda\sum_i |\beta_i| .
\]
In sklearn's \texttt{Lasso} (and statsmodels' \texttt{fit\_regularized}) this parameter
is named \texttt{alpha}, so each $\lambda$ is passed in as \texttt{alpha=lam}. In the
lab's elastic-net notation, pure lasso fixes the mixing weight \texttt{L1\_wt} $= 1$, which
leaves the penalty strength as the only parameter to tune.

We use sklearn's \texttt{Lasso}, which minimizes exactly the loss above (the same
$\frac{1}{2n}$ scaling as statsmodels), so the values are directly comparable.
statsmodels' lasso solver runs coordinate descent in pure Python and was too slow with
3872 features. At the smallest $\lambda$ values the solver does not fully converge
(\texttt{ConvergenceWarning}). Those models overfit and are not candidates for the
best $\lambda$, so the warnings are suppressed.

\subsubsection*{Coarse sweep}

\begin{lstlisting}
import warnings
from sklearn.exceptions import ConvergenceWarning
from sklearn.linear_model import Lasso

warnings.filterwarnings("ignore", category=ConvergenceWarning)

lambdas = np.logspace(-6, 3, 10)
r_sqds = []
n_nonzero = []

for lam in lambdas:
    lasso = Lasso(alpha=lam, max_iter=10_000)
    lasso.fit(x_train, y_train)
    r_sqds.append(r2_score(y_test, lasso.predict(x_test)))
    n_nonzero.append(np.sum(lasso.coef_ != 0))     # pixels kept

plt.semilogx(lambdas, r_sqds, "o-")
plt.xlabel("$\\lambda$ (lasso penalty strength)"); plt.ylabel("test $r^2$")
plt.show()
\end{lstlisting}

\begin{center}
\begin{tabular}{lccccc}
\toprule
$\lambda$ & $10^{-6}$ & $10^{-5}$ & $10^{-4}$ & $10^{-3}$ & $\ge 10^{-2}$ \\
\midrule
test $r^2$    & $0.004$ & $0.200$ & $0.366$ & $-0.048$ & $-0.048$ \\
pixels kept   & 496     & 377     & 77      & 0        & 0 \\
\bottomrule
\end{tabular}
\end{center}

\begin{figure}[h]
  \centering
  \includegraphics[width=0.6\textwidth]{figures/lasso_coarse.png}
  \caption{Lasso: test $r^2$ against $\lambda$, coarse sweep (log scale).}
\end{figure}

\subsubsection*{Zooming in}

The coarse sweep puts the peak near $\lambda = 10^{-4}$. A first zoom over 30 values
from $10^{-5}$ to $10^{-3}$ placed the best value at $\lambda \approx 4.4 \times 10^{-5}$.
Its grid points are $3.4 \times 10^{-5}$ apart, which is wide compared with $\lambda$ itself
near the peak, so a second zoom over 13 values from $2 \times 10^{-5}$ to $8 \times 10^{-5}$
was used to pin it down. Both zooms use the same loop as above with
\texttt{np.linspace}:

\begin{lstlisting}
lambdas = np.linspace(0.00001, 0.001, 30)     # zoom 1
lambdas = np.linspace(0.00002, 0.00008, 13)   # zoom 2
\end{lstlisting}

\begin{figure}[h]
  \centering
  \includegraphics[width=0.48\textwidth]{figures/lasso_zoom1.png}
  \hfill
  \includegraphics[width=0.48\textwidth]{figures/lasso_zoom2.png}
  \caption{Lasso: test $r^2$ against $\lambda$. Left: zoom 1 ($10^{-5}$ to $10^{-3}$).
  Right: zoom 2 ($2\times10^{-5}$ to $8\times10^{-5}$).}
\end{figure}

\begin{center}
\begin{tabular}{lccccccc}
\toprule
$\lambda$ ($\times 10^{-5}$) & 2.0 & 3.0 & 4.0 & 5.0 & \textbf{5.5} & 6.0 & 7.0 \\
\midrule
test $r^2$  & $0.320$ & $0.356$ & $0.378$ & $0.389$ & $\mathbf{0.389}$ & $0.389$ & $0.381$ \\
pixels kept & 300     & 249     & 202     & 168     & \textbf{148}     & 135     & 105 \\
\bottomrule
\end{tabular}
\end{center}

\textbf{Answer:} The best lasso penalty is
$\boldsymbol{\lambda \approx 5.5 \times 10^{-5}}$, with test $r^2 \approx 0.389$. At
this $\lambda$, lasso keeps only \textbf{148 of the 3872 pixels} (about 4\%) and sets the rest
exactly to zero.

\textbf{Justification.}
\begin{itemize}
  \item \textbf{$\lambda$ too small:} the penalty is too weak, so lasso keeps hundreds of
  pixels and overfits the 487 training images. At $\lambda = 10^{-6}$ it keeps 496 pixels
  and the test $r^2$ is only $0.004$. At $\lambda = 10^{-5}$ it keeps 377 pixels with test
  $r^2 = 0.200$. The solver also fails to converge here, a sign that too many correlated
  pixels are competing to explain the same signal.
  \item \textbf{$\lambda$ too large:} the penalty sets more and more coefficients to exactly
  zero. By $\lambda \approx 5 \times 10^{-4}$ only about 10 pixels remain and test $r^2$ has
  fallen to $0.03$. For $\lambda \ge 8 \times 10^{-4}$ every pixel is dropped, the model
  predicts the mean \texttt{nWBV} for every image, and test $r^2 = -0.048$.
  \item \textbf{$\lambda \approx 5.5 \times 10^{-5}$} balances the two. It discards the pixels
  that do not help and keeps a small set that predicts brain volume.
\end{itemize}
As with ridge, the peak is flat: $\lambda$ from $5 \times 10^{-5}$ to $6 \times 10^{-5}$
all give test $r^2 \approx 0.389$, keeping between 135 and 168 pixels.

\subsubsection*{Comparison with ridge}

\begin{center}
\begin{tabular}{lccc}
\toprule
Method & Best penalty & Test $r^2$ & Pixels used \\
\midrule
Ridge & $\alpha \approx 0.009$             & $0.398$ & 3872 (all, shrunk) \\
Lasso & $\lambda \approx 5.5\times10^{-5}$ & $0.389$ & 148 \\
\bottomrule
\end{tabular}
\end{center}

Lasso nearly matches ridge's accuracy with a far simpler model: about 4\% of the
pixels. Ridge's small edge fits how brain volume appears in an MRI slice. It shows up
as the overall amount of tissue against dark space across the whole image rather than
in a few specific pixels. Spreading small weights over every pixel (ridge) suits this
kind of diffuse signal slightly better than selecting a few pixels (lasso).

\end{document}
